\section{Introduction}

Molecular dynamics (MD) is the workhorse for connecting microscopic interactions to the
macroscopic behaviour of liquids, solids, and biological assemblies. Its central limitation is
time: the 1--2 fs integration step required for stable dynamics restricts routine simulations to
nanoseconds, while many phenomena of interest---diffusion, nucleation, conformational transitions,
and slow relaxation---unfold over microseconds to milliseconds. Machine-learned interatomic
potentials have reduced the cost of the force evaluation, but they retain the sequential
integrator and therefore the same step-count bottleneck.

A more radical route is to learn the propagator itself. Instead of evaluating forces and
integrating Newton's equations, a direct propagator maps a state at time \(t\) to a state at
\(t+\Delta t\) with \(\Delta t\) several times the stable integration step. Because the map is
force-free, it can in principle be evaluated at a fraction of the cost of an MD step and can be
rolled out autoregressively. TrajCast demonstrated that this idea is viable for liquid water: an
equivariant message-passing network trained on 5 fs pairs reproduces the oxygen--oxygen radial
distribution function (RDF), the vibrational density of states (VDOS), and the diffusion
coefficient under a global thermostat \cite{trajcast}. The result is an important existence proof:
faithful long-time water physics is possible without forces.

What the existence proof does not provide is a guarantee. A direct propagator is trained on
one-step pairs, usually by minimizing a displacement or velocity error. The quantity that matters
at long times, however, is the stationary measure of the composed map---the distribution of
positions and velocities that the autoregressive rollout visits after transients have decayed.
One-step accuracy constrains a conditional mean; it says little about the invariant measure of a
nonlinear, stochastic, or thermostatted iteration. The gap is familiar in spirit from generative
modelling, where a low reconstruction error does not imply a correct sample distribution, but it
has not been diagnosed systematically for force-free MD propagators.

Here we show that the gap is real and that it has a specific structure. Using water as a test
system, we compare a published TrajCast checkpoint with our own direct propagator. The published
model has a one-step position error about fifteen times larger than ours on held-out transitions,
yet it reproduces the canonical mode partition, RDF, and VDOS, whereas our pointwise-better model
develops a large mode imbalance and fails the same 100 ps NVT screen. The ranking by one-step
error is therefore not just noisy---it can be inverted. A second diagnostic removes another
plausible explanation: 25 ps NVE energy drift does not order the models either. The published
model heats strongly in NVE and still passes NVT; our most NVE-stable model fails NVT. The
relevant object is the composed propagator plus thermostat, and the relevant failure is a bias in
the mode-resolved stationary measure, not a one-step error or an energy-integration error. This
is an existence proof from two models, not a controlled correlation: the checkpoints differ in
architecture and training data, so the inversion establishes that one-step accuracy is neither
necessary nor sufficient, not that it is generally anti-correlated with ensemble fidelity.

Existing remedies attack different parts of this problem. Global stochastic thermostats, such as
CSVR, control the total kinetic energy but not its distribution among molecular translation,
rotation, and internal motion \cite{csvr}. Structure-preserving learned maps enforce symplectic or
time-reversible structure and improve energy conservation and equipartition \cite{bigi_action},
but they are deterministic and target mechanical structure rather than the canonical NVT measure.
Learned velocity updates recover stability by injecting an LSTM prediction periodically
\cite{mamd}, but the update is a black box and is not tied to the equilibrium mode partition.
Generalized Langevin equation (GLE) thermostats provide memory kernels and colored noise with the
correct fluctuation--dissipation relation \cite{ceriotti_prl,ceriotti_jctc}, and
frequency-resolved thermostats can control selected modes \cite{zpe_thermostat}, but these are
thermostats for an MD integrator, not readouts for a frozen force-free propagator. Stochastic and
generative propagators offer another route---Timewarp accelerates dynamics with learned
time-coarsening \cite{timewarp}, and generative models such as MDGen sample trajectory
distributions \cite{mdgen}---at the cost of MCMC, an energy function, or a substantially larger
generative model.

We take a different decomposition. A frozen force-free backbone is responsible for structure: it
supplies the displacement \(\Delta q_t\). A lightweight, mode-resolved generalized-Langevin
readout (GLR) is responsible for thermodynamics: it supplies the velocity that enters the next
backbone call and the kinetic diagnostics. The readout decomposes the velocity into molecular
translation, rotation, and internal modes; fits a per-mode AR(\(p\)) memory kernel to the
reference velocity autocorrelation by Yule--Walker; injects fluctuation--dissipation-consistent
noise with the prediction-error variance; and rescales each mode or group to its canonical kinetic-energy target
energy \(N_a k_B T/2\) at every step. The partition is therefore enforced by construction, while
the memory kernel matches the multi-lag dynamics that a single gain cannot capture. The backbone
is not retrained, and the readout is fitted from a short reference NVT trajectory.

We apply the GLR to SPC/Fw water \cite{spcfw} and compare it with the published TrajCast
checkpoint, our 250k direct propagator, an AR(1) readout, and several negative controls. The GLR
is stable for 100 ps from two independent initial conditions, is initial-condition independent in
its partition, ACF, RDF, and VDOS, and reproduces the O--O, O--H, and H--H RDFs and the O VDOS
more closely than the published checkpoint in our analysis. The 2.5 ns production run tests the
diffusion coefficient---the last long-time observable---and the ablations isolate the roles of
the memory order, the FD noise, the variance rescale, and the displacement coupling. The result is
a practical route, demonstrated here on water, for making a force-free propagator
thermodynamically consistent without retraining it, and a diagnostic that explains when and why
such a propagator will fail.

\subsection*{Contributions}

\begin{enumerate}
\item A diagnostic showing that pointwise one-step accuracy is neither necessary nor sufficient
  for long-time equilibrium, and that NVE energy stability is not the discriminating variable.
\item Evidence that the failure is a mode-resolved stationary-measure bias of the composed
  propagator plus thermostat.
\item A mode-resolved AR(\(p\)) FD readout with exact group-resolved variance rescaling that
  guarantees the canonical mean kinetic partition (and the marginal with the stochastic target)
  and matches multi-lag mode ACFs on a frozen force-free
  propagator.
\item Water results at or beyond the published TrajCast checkpoint on RDF and VDOS, with diffusion
  and initial-condition independence tested at 2.5 ns, plus ablations that isolate each component.
\end{enumerate}
